% ---------------------------------------------------------
% TABELA EMENTA 40: Gestão de Marketing II
% ---------------------------------------------------------
\begin{xltabular}{\linewidth}{|X|p{2.5cm}|p{2.5cm}|}
\hline
\multirow{3}{=}{\textcolor{ifscgreen}{\textbf{Unidade Curricular:}}\\[3pt]\textbf{\large Gestão de Marketing II}} & \multicolumn{2}{p{\dimexpr5cm+2\tabcolsep+\arrayrulewidth\relax}|}{\textcolor{ifscgreen}{\textbf{Semestre:}} \textbf{3º (Ano 2)}} \\ \cline{2-3}
 & \textcolor{ifscgreen}{\textbf{CH EaD*:}} & \textcolor{ifscgreen}{\textbf{CH Total*:}} \\ \cline{2-3}
 & \textbf{00 h} & \textbf{40 h} \\ \hline
\endhead
\multicolumn{3}{|p{\dimexpr\linewidth-2\tabcolsep-2\arrayrulewidth\relax}|}{\textcolor{ifscgreen}{\textbf{Objetivos:}}} \\ \hline
\multicolumn{3}{|p{\dimexpr\linewidth-2\tabcolsep-2\arrayrulewidth\relax}|}{
\begin{itemize}[leftmargin=*,noitemsep,topsep=2pt]
 \item \textit{[Objetivos de aprendizagem pendentes de preenchimento pelo docente responsável]}
\end{itemize}
} \\ \hline
\multicolumn{3}{|p{\dimexpr\linewidth-2\tabcolsep-2\arrayrulewidth\relax}|}{\textcolor{ifscgreen}{\textbf{Conteúdos:}}} \\ \hline
\multicolumn{3}{|p{\dimexpr\linewidth-2\tabcolsep-2\arrayrulewidth\relax}|}{
\begin{itemize}[leftmargin=*,noitemsep,topsep=2pt]
 \item \textit{[Conteúdo programático pendente de preenchimento pelo docente responsável]}
\end{itemize}
} \\ \hline
\multicolumn{3}{|p{\dimexpr\linewidth-2\tabcolsep-2\arrayrulewidth\relax}|}{\textcolor{ifscgreen}{\textbf{Estratégias de Ensino e Aprendizagem:}}} \\ \hline
\multicolumn{3}{|p{\dimexpr\linewidth-2\tabcolsep-2\arrayrulewidth\relax}|}{\textit{[Metodologia de ensino e avaliação pendente de preenchimento pelo docente responsável]}} \\ \hline
\multicolumn{3}{|p{\dimexpr\linewidth-2\tabcolsep-2\arrayrulewidth\relax}|}{\textcolor{ifscgreen}{\textbf{Bibliografia Básica:}}} \\ \hline
\multicolumn{3}{|p{\dimexpr\linewidth-2\tabcolsep-2\arrayrulewidth\relax}|}{1. \textit{[1º título da Bibliografia Básica pendente de indicação docente e validação no Parecer da Biblioteca do Câmpus]}
2. \textit{[2º título da Bibliografia Básica pendente de indicação docente e validação no Parecer da Biblioteca do Câmpus]}} \\ \hline
\multicolumn{3}{|p{\dimexpr\linewidth-2\tabcolsep-2\arrayrulewidth\relax}|}{\textcolor{ifscgreen}{\textbf{Bibliografia Complementar:}}} \\ \hline
\multicolumn{3}{|p{\dimexpr\linewidth-2\tabcolsep-2\arrayrulewidth\relax}|}{1. \textit{[1º título da Bibliografia Complementar pendente de indicação docente e validação no Parecer da Biblioteca do Câmpus]}
2. \textit{[2º título da Bibliografia Complementar pendente de indicação docente e validação no Parecer da Biblioteca do Câmpus]}
3. \textit{[3º título da Bibliografia Complementar pendente de indicação docente e validação no Parecer da Biblioteca do Câmpus]}} \\ \hline
\end{xltabular}

\vspace{0.4cm}


% ---------------------------------------------------------
% TABELA EMENTA 41: Gestão de Operações e Qualidade
% ---------------------------------------------------------
\begin{xltabular}{\linewidth}{|X|p{2.5cm}|p{2.5cm}|}
\hline
\multirow{3}{=}{\textcolor{ifscgreen}{\textbf{Unidade Curricular:}}\\[3pt]\textbf{\large Gestão de Operações e Qualidade}} & \multicolumn{2}{p{\dimexpr5cm+2\tabcolsep+\arrayrulewidth\relax}|}{\textcolor{ifscgreen}{\textbf{Semestre:}} \textbf{4º e 5º (Ano 2/3)}} \\ \cline{2-3}
 & \textcolor{ifscgreen}{\textbf{CH EaD*:}} & \textcolor{ifscgreen}{\textbf{CH Total*:}} \\ \cline{2-3}
 & \textbf{00 h} & \textbf{80 h} \\ \hline
\endhead
\multicolumn{3}{|p{\dimexpr\linewidth-2\tabcolsep-2\arrayrulewidth\relax}|}{\textcolor{ifscgreen}{\textbf{Objetivos:}}} \\ \hline
\multicolumn{3}{|p{\dimexpr\linewidth-2\tabcolsep-2\arrayrulewidth\relax}|}{
\begin{itemize}[leftmargin=*,noitemsep,topsep=2pt]
 \item Auxiliar no planejamento, na organização, na execução e no controle das operações relacionadas à produção de bens, à prestação de serviços, aos suprimentos, à logística e à gestão da qualidade.
 \item Executar rotinas administrativas relacionadas às operações organizacionais, aplicando princípios da gestão de operações e da qualidade.
 \item Identificar, sistematizar e interpretar informações do contexto organizacional para subsidiar a tomada de decisão e a proposição de melhorias nas operações.
\end{itemize}
} \\ \hline
\multicolumn{3}{|p{\dimexpr\linewidth-2\tabcolsep-2\arrayrulewidth\relax}|}{\textcolor{ifscgreen}{\textbf{Conteúdos:}}} \\ \hline
\multicolumn{3}{|p{\dimexpr\linewidth-2\tabcolsep-2\arrayrulewidth\relax}|}{
\begin{itemize}[leftmargin=*,noitemsep,topsep=2pt]
 \item Conceitos e princípios da gestão de operações; operações de produção de bens e de prestação de serviços; planejamento e controle de operações.
 \item Previsão de demanda e planejamento da capacidade produtiva.
 \item Gestão de compras, suprimentos e fornecedores.
 \item Gestão de estoques.
 \item Gestão de logística e cadeia de suprimentos; logística reversa.
 \item Conceitos e princípios da gestão da qualidade; evolução histórica e abordagens contemporâneas.
 \item Ferramentas da gestão da qualidade (fluxograma, diagrama de causa e efeito, Pareto, 5W2H, ciclo PDCA).
 \item Sistemas de gestão da qualidade.
 \item Operações, qualidade e sustentabilidade nas organizações.
\end{itemize}
} \\ \hline
\multicolumn{3}{|p{\dimexpr\linewidth-2\tabcolsep-2\arrayrulewidth\relax}|}{\textcolor{ifscgreen}{\textbf{Estratégias de Ensino e Aprendizagem:}}} \\ \hline
\multicolumn{3}{|p{\dimexpr\linewidth-2\tabcolsep-2\arrayrulewidth\relax}|}{- **Metodologia:** Aulas expositivas dialogadas, com enfoque teórico-prático, incluindo atividades de simulação e dinâmicas envolvendo planejamento e controle de operações, estudos de casos aplicados às organizações e aos arranjos produtivos, análise de documentos e indicadores de desempenho, e interações com empreendimentos e profissionais da área. Utiliza sala de aula, laboratório de informática, sala multidisciplinar e espaços de organizações externas. Articula-se com Introdução à Administração e Organização e Processos, além de Língua Portuguesa, Matemática, Sociologia, Sociedade e Trabalho, História e Geografia.
- **Avaliação:** Contínua, processual e formativa, contemplando produções escritas, apresentações orais, resolução de estudos de caso, atividades práticas, relatórios, registros reflexivos e participação nas atividades propostas, de forma individual ou em grupo.} \\ \hline
\multicolumn{3}{|p{\dimexpr\linewidth-2\tabcolsep-2\arrayrulewidth\relax}|}{\textcolor{ifscgreen}{\textbf{Bibliografia Básica:}}} \\ \hline
\multicolumn{3}{|p{\dimexpr\linewidth-2\tabcolsep-2\arrayrulewidth\relax}|}{1. BALLOU, Ronald H. \textbf{Gerenciamento da cadeia de suprimentos: planejamento, organização e logística empresarial}. 4. ed. Porto Alegre: Bookman, 2001.
2. CORRÊA, Carlos Alberto; CORRÊA, Henrique Luiz. \textbf{Administração de produção e operações: manufatura e serviços, uma abordagem estratégica}. São Paulo: Atlas, 2006.
3. DIAS, Marco Aurélio P. \textbf{Administração de materiais: uma abordagem logística}. 5. ed. São Paulo: Atlas, 2010.
4. GIANESI, Irineu G. N.; CORRÊA, Henrique Luiz. \textbf{Administração estratégica de serviços: operações para a satisfação do cliente}. São Paulo: Atlas, 2013.
5. PALADINI, Edson P. \textbf{Gestão estratégica da qualidade: princípios, métodos e processos}. 2. ed. São Paulo: Atlas, 2009.} \\ \hline
\multicolumn{3}{|p{\dimexpr\linewidth-2\tabcolsep-2\arrayrulewidth\relax}|}{\textcolor{ifscgreen}{\textbf{Bibliografia Complementar:}}} \\ \hline
\multicolumn{3}{|p{\dimexpr\linewidth-2\tabcolsep-2\arrayrulewidth\relax}|}{1. ARBACHE, F. S. et al. \textbf{Gestão de logística, distribuição e trade marketing}. Rio de Janeiro: FGV, 2004.
2. LEITE, Paulo Roberto. \textbf{Logística reversa: sustentabilidade e competitividade}. 3. ed. São Paulo: Saraiva, 2017.
3. MARTINS, Petrônio G.; LAUGENI, Fernando P. \textbf{Administração da produção}. São Paulo: Saraiva, 2006.
4. MARSHALL JUNIOR, Isnard. \textbf{Gestão da qualidade}. 10. ed. Rio de Janeiro: Ed. da FGV, 2010.
5. MOREIRA, Daniel Augusto. \textbf{Administração da produção e operações}. São Paulo: Cengage Learning, 2009.
6. NOGUEIRA, José Francisco. \textbf{Gestão estratégica de serviços: teoria e prática}. São Paulo: Atlas, 2008.} \\ \hline
\end{xltabular}

\vspace{0.4cm}

